\documentclass[UTF8,a4paper,11pt]{ctexart}
\usepackage[a4paper,margin=2.3cm]{geometry}
\usepackage{amsmath,amssymb,mathtools}
\usepackage{enumitem}
\usepackage[hidelinks]{hyperref}
\usepackage{fancyhdr}
\setlist[itemize]{leftmargin=2.2em,itemsep=0.25em,topsep=0.35em}
\setlength{\parindent}{2em}
\setlength{\parskip}{0.35em}
\pagestyle{fancy}\fancyhf{}
\fancyhead[L]{离散数学与结构（I）}\fancyhead[R]{Lec 08}
\fancyfoot[C]{\thepage}\setlength{\headheight}{14pt}
\title{共轭、类方程与 Sylow 定理}
\author{zhangjunshuo}
\date{2026 年 10 月}
\begin{document}
\maketitle
\begin{quote}
本讲根据 2026 年 10 月 8 日的两份课堂转写、手写原稿和教材整理。转写中的口语识别错误已按手写稿与 Dummit--Foote《Abstract Algebra》（[DF]）和《离散数学与结构》（[DS]）校正；为使证明可读，个别课堂只给出结论的地方补写了中间步骤，并明确标注整理补全。
\end{quote}
\tableofcontents
\clearpage

\section*{记号速查}
\begin{itemize}
\item $G,H,K$：群或子群；$e$：单位元；$|G|$：群的阶；$H\le G$：子群；$H\trianglelefteq G$：$H$ 正规于 $G$。
\item $\langle x\rangle$：由 $x$ 生成的循环子群；$C_n$：阶 $n$ 的循环群；$S_n$：$n$ 个字母的对称群。
\item $Z(G)$：中心；$C_G(x)$：元素的中心化子；$N_G(H)$：子群的正规化子。
\item $G\curvearrowright X$：$G$ 作用在 $X$ 上；$G\cdot x$：轨道；$G_x$：稳定子；$X^g$：固定点集。
\item $|G|=p^a m$ 且 $p\nmid m$ 时，Sylow $p$-子群的阶为 $p^a$，全体记为 $\operatorname{Syl}_p(G)$，个数记为 $n_p$。
\item $[G:H]=|G|/|H|$：指数；$[x,y]=xyx^{-1}y^{-1}$：交换子。
\end{itemize}
\clearpage

\section{组合列、单群与有限群分解}


本节先回答一个结构问题：阿贝尔群能拆成循环群直积；一般有限群通常不能直接拆成直积，但可以沿着正规子群链逐层拆成\textbf{单群商}。这是一种“逐层观察”的分解，不应误读为把原群直接写成所有商群的直积。

\textbf{本节用到的记号，先一次交代清楚}（更早的讲次已逐个定义过，这里只复述结论，方便单独阅读本节）：

\begin{itemize}
\item $e$：群 $G$ 的单位元；$\{e\}$ 是平凡子群。
\item $|G|$：群 $G$ 的阶，即 $G$ 中元素的个数。
\item $H\le G$：$H$ 是 $G$ 的子群；$H\trianglelefteq G$：$H$ 是 $G$ 的正规子群（对任意 $g\in G$ 有 $gHg^{-1}=H$）。
\item $G/N$：$N\trianglelefteq G$ 时的商群，元素是陪集 $gN$；$[G:H]=|G|/|H|$ 是 $H$ 在 $G$ 中的指数。
\item $\cong$：同构；$G\cong H$ 表示两个群有完全相同的乘法结构。
\item $C_n$：$n$ 阶循环群，即由一个元素生成的 $n$ 阶群（例如 $C_3=\{e,g,g^2\}$，$g^3=e$）。
\item $S_n$：$n$ 个字母上的对称群，即 $\{1,\dots,n\}$ 到自身的全体双射在复合下构成的群，$|S_n|=n!$。置换写成轮换的乘积，例如 $(1\,2\,3)$ 表示 $1\mapsto2\mapsto3\mapsto1$，$(1\,2)$ 表示交换 $1,2$。这套记号与置换的复合约定见 [Lec 06 §1.3 对称群](../lec06/sections/03-permutations.md)。
\item $A_n\le S_n$：偶置换组成的交错群；本节的 $A_3=\{e,(1\,2\,3),(1\,3\,2)\}$ 是 $S_3$ 中由三循环和恒等元组成的子群。
\item $\times$：群的直积（按分量运算）。
\end{itemize}


\subsection{正规链与单群}


设 $G$ 是群。一条正规链（normal chain）是

\[
\{e\}=G_0\trianglelefteq G_1\trianglelefteq\cdots\trianglelefteq G_r=G,
\]

其中每个 $G_i$ 只要求正规于它紧邻的上一级 $G_{i+1}$（即 $G_i\trianglelefteq G_{i+1}$，不是 $G_i\trianglelefteq G$）。一般不能推出 $G_i\trianglelefteq G$；这正是读组合列时最容易混淆的地方。每一层产生一个商群 $G_{i+1}/G_i$。

一个非平凡群 $S$ 称为\textbf{单群}（simple group），如果它没有非平凡真正规子群：

\[
N\trianglelefteq S\quad\Longrightarrow\quad N=\{e\}\ \text{或}\ N=S.
\]

这里“非平凡”很重要；平凡群通常不作为单群因子。单群是正规分解中已经不能再沿正规子群继续切开的“最小砖块”。

如果一条正规链的每个因子 $G_{i+1}/G_i$ 都是单群，就称它为 \textbf{composition series（组合列）}，这些单群商称为组合因子。

\subsection{为什么有限群总能得到组合列}


\textbf{命题（存在性）。} 每个有限群都有组合列。

可以按群的阶 $|G|$（即元素个数）归纳。若 $G$ 本身单，则

\[
\{e\}\trianglelefteq G
\]

已经是一条组合列。若 $G$ 非单，取一个非平凡真正规子群 $N\trianglelefteq G$。归纳假设给出 $N$ 的组合列；同时 $|G/N|<|G|$，所以 $G/N$ 也有组合列：

\[
\{eN\}=\overline G_0\trianglelefteq\overline G_1\trianglelefteq\cdots\trianglelefteq\overline G_s=G/N.
\]

把商群中的链取原像，就得到 $G$ 中从 $N$ 到 $G$ 的一段链。具体地，令商映射 $\pi:G\to G/N$，$\pi(g)=gN$（这里的 $\overline G_j$ 表示 $\overline G_j\subseteq G/N$，上划线只是提示它住在商群里），则

\[
H_j=\pi^{-1}(\overline G_j)
\]

满足 $N=H_0\trianglelefteq H_1\trianglelefteq\cdots\trianglelefteq H_s=G$，并且由对应定理

\[
H_{j+1}/H_j\cong \overline G_{j+1}/\overline G_j
\]

都是单群。把 $N$ 内的链和这段链接起来即可。有限性保证每次选择真子群都会严格减小群的阶，归纳必然终止。

另一种课堂上常用的说法是：若某一层商群还不是单群，就在该商群中取一个非平凡真正规子群，再把它取原像插回原链；有限群不可能无限插入。

\subsection{例：$S_3$ 的组合列}


在对称群 $S_3$（$\{1,2,3\}$ 上的置换群，$|S_3|=6$）中，三循环和恒等元组成的子群

\[
A_3=\{e,(123),(132)\}
\]

是指数 $2$ 的正规子群。因此

\[
\{e\}\trianglelefteq A_3\trianglelefteq S_3
\]

是一条组合列，因为

\[
A_3/\{e\}\cong C_3,
\qquad S_3/A_3\cong C_2
\]

都是素数阶循环群，因而是单群。

这条链表达的是“先看到一个 $C_3$ 层，再看到一个 $C_2$ 层”。它\textbf{不}表示 $S_3\cong C_3\times C_2$：后者是阿贝尔群，而 $S_3$ 非阿贝尔。组合列的商群只是逐层记录正规结构，不能自动恢复为直积。

\subsection{组合因子的唯一性与课堂范围}


Jordan--Hölder 定理说，有限群的任意两条组合列得到的组合因子，在忽略排列顺序和同构后相同。它类似整数分解的“素因子多重集合唯一”，但这里的因子是单群，不一定是阿贝尔群。

有限单群分类定理进一步描述所有有限单群的类型；证明由大量结果组成，并不是本节或本课程需要掌握的证明。当前最常用的单群例子是素数阶循环群 $C_p$，以及后续可能遇到的非阿贝尔单群。

\subsection{读题时的三个检查}


\begin{itemize}
\item 看到 $G_i\trianglelefteq G_{i+1}$，先确认正规性只针对相邻两项；除非另有说明，不要写成 $G_i\trianglelefteq G$。
\item 看到 $G_{i+1}/G_i$，先检查 $G_i$ 是否正规于 $G_{i+1}$，否则商群没有定义。
\item “因子是单群”描述的是商群的内部结构；“$G$ 是若干群的直积”是另一种更强的结论，需要额外的正规性、交平凡和乘积条件。
\end{itemize}


\subsection{参考对照}


\begin{itemize}
\item [DF] 关于正规列、组合列和 Jordan--Hölder 定理的章节。
\item [DS] 关于正规子群、商群和对应定理的章节；本节的原像构造直接使用对应定理。
\end{itemize}

\section{Burnside 引理与对称群中的共轭类}


本节先用一次双计数得到 Burnside 引理，再把共轭作用应用于 $S_n$。前者数轨道，后者给出一个可操作的共轭类判别：在对称群里，共轭类由不交循环的长度决定。

Burnside 引理本身对任意群作用都成立，不需要“群作用在自己身上”。下一节要做的共轭作用只是它的一个特例；先把“平均固定点数”这件工具单独学会，再用到共轭类上会更清楚。

\textbf{本节用到的记号，先一次交代清楚}（群作用的一般理论在 [Lec 07 §4–5](../lec07/sections/04-group-actions.md) 已建立，这里只复述本节要用的部分）：

\begin{itemize}
\item 群 $G$ \textbf{作用}在集合 $X$ 上，记作 $G\curvearrowright X$，是指映射 $G\times X\to X$，$(g,x)\mapsto g\cdot x$，满足 $e\cdot x=x$ 与 $(gh)\cdot x=g\cdot(h\cdot x)$。
\item $x$ 的\textbf{轨道}为 $G\cdot x=\{g\cdot x:g\in G\}$；\textbf{稳定子}为 $G_x=\{g\in G:g\cdot x=x\}$，它是 $G$ 的子群。
\item $X/G$ 表示全体轨道组成的集合，因此 $|X/G|$ 就是\textbf{轨道的个数}。
\item $[G:H]=|G|/|H|$ 是子群 $H\le G$ 的\textbf{指数}；轨道--稳定子定理说 $|G\cdot x|=[G:G_x]$。
\item $|S|$ 表示集合 $S$ 的元素个数；$\#S$ 是同一个意思，本讲在集合计数式子里偶尔使用。
\item 对称群 $S_n$ 与轮换记号 $(1\,2\,3)$ 见 [Lec 06 §1.3](../lec06/sections/03-permutations.md)。\textbf{共轭作用}指 $g\cdot x=gxg^{-1}$，它的轨道叫\textbf{共轭类}；下一节会正式展开。
\item $C_n$ 是 $n$ 阶循环群；本节出现的 $C_4=\langle r\mid r^4=e\rangle$ 由 $90^\circ$ 旋转 $r$ 生成，$D_8$ 是正方形的全部 8 个对称（旋转与翻转）构成的群。
\end{itemize}


\subsection{Burnside 引理}


设有限群 $G$ 作用在有限集合 $X$ 上。对 $g\in G$，记固定点集为

\[
X^g=\operatorname{Fix}_X(g)=\{x\in X:g\cdot x=x\}.
\]

\textbf{Burnside 引理（Cauchy--Frobenius）。} $G$ 在 $X$ 上的轨道数为

\[
|X/G|=\frac{1}{|G|}\sum_{g\in G}|X^g|.
\]

注意 $X^g$ 是固定一个群元素 $g$ 后留下的点集；它和固定一个点 $x$ 的稳定子 $G_x$ 是两个不同对象。

\subsubsection{双计数证明}


考虑集合

\[
\Omega=\{(g,x)\in G\times X:g\cdot x=x\}.
\]

先固定 $g$。允许的 $x$ 正好构成 $X^g$，所以

\[
|\Omega|=\sum_{g\in G}|X^g|.
\]

再固定一个轨道 $O=G\cdot x$。轨道中每个 $y$ 的稳定子都与 $G_x$ 共轭，故大小相同。于是

\[
\#\{(g,y)\in G\times O:g\cdot y=y\}
=|O|\,|G_x|
=[G:G_x]|G_x|=|G|.
\]

每条轨道贡献恰好 $|G|$ 个点对；若轨道总数为 $r$，则 $|\Omega|=r|G|$。与前一个计数相等并除以 $|G|$，即得公式。

\subsubsection{例：正方形顶点的二染色（只允许旋转）}


令 $G=C_4=\langle r\mid r^4=e\rangle$ 旋转正方形，令 $X=\{0,1\}^4$ 表示四个顶点的黑白染色，共有 $16$ 种。我们把旋转后相同的染色视为同一种。

\begin{itemize}
\item 恒等旋转固定全部 $16$ 种染色；
\item $90^\circ$ 和 $270^\circ$ 旋转要求四个顶点颜色相同，各固定 $2$ 种；
\item $180^\circ$ 旋转要求相对顶点同色，固定 $2^2=4$ 种。
\end{itemize}


因此不同旋转轨道数为

\[
\frac{16+2+4+2}{4}=6.
\]

\subsubsection{同一个例子的完整版：把反射也算进去}


如果认为翻转后重合的染色也是同一种，群就要换成正方形的全部 $8$ 个对称

\[
D_8=\langle r,s\mid r^4=s^2=e,\ srs^{-1}=r^{-1}\rangle,
\]

其中 $r$ 是 $90^\circ$ 旋转，$s$ 是固定顶点 $1,3$ 的反射。顶点按顺时针编号为 $1,2,3,4$，则 $D_8$ 的 $8$ 个元素是 $r^k$ 与 $sr^k$（$k=0,1,2,3$）。

数固定染色数有一个统一办法：把置换写成不交循环，\textbf{每个循环上的顶点必须同色}，所以固定染色数等于 $2^{\text{循环个数}}$。逐个元素数出来：


\begin{center}
\begin{tabular}{lll}
群元素 & 循环分解 & 固定染色数\\ \hline
$e$ & $(1)(2)(3)(4)$ & $2^4=16$\\
$r,\ r^3$ & $(1\,2\,3\,4)$ & $2^1=2$\\
$r^2$ & $(1\,3)(2\,4)$ & $2^2=4$\\
$s$（轴过顶点 $1,3$） & $(1)(3)(2\,4)$ & $2^3=8$\\
$sr$（轴过顶点 $2,4$） & $(2)(4)(1\,3)$ & $2^3=8$\\
$sr^2$（轴过对边 $12$、$34$ 的中点） & $(1\,2)(3\,4)$ & $2^2=4$\\
$sr^3$（轴过对边 $23$、$41$ 的中点） & $(1\,4)(2\,3)$ & $2^2=4$
\end{tabular}
\end{center}


Burnside 引理给出

\[
\frac{16+2+4+2+8+8+4+4}{8}=\frac{48}{8}=6.
\]

和只用旋转的 $C_4$ 相比，轨道数仍然是 $6$。\textbf{而且这里连轨道划分都完全一样}，原因是查表就能发现的巧合：每个反射作用在染色上的效果恰好等于某个旋转。


\begin{center}
\begin{tabular}{lll}
反射元素 & 它在染色上的作用 & 等于哪个旋转\\ \hline
$s$ & 交换顶点 $2,4$，固定 $1,3$ & $r^2$\\
$sr$ & 交换顶点 $1,3$，固定 $2,4$ & $r^2$\\
$sr^2$ & 交换顶点 $1,2$ 与 $3,4$ & $r$ 或 $r^3$\\
$sr^3$ & 交换顶点 $1,4$ 与 $2,3$ & $r$ 或 $r^3$
\end{tabular}
\end{center}


这四次“巧合”不是偶然：$4$ 个顶点的二染色恰好被对边关系决定，而每个反射都把顶点配成固定的对子，于是它对染色的作用只是一个坐标置换，正好落在旋转已经产生的置换里。所以两个群给出同一个划分，轨道大小都是

\[
4,\ 4,\ 4,\ 2,\ 1,\ 1,
\]

例如 $0011$（顶点 $1,2$ 黑，$3,4$ 白）在两个群下的轨道都是

\[
\{0011,\ 0110,\ 1001,\ 1100\}.
\]

读这个例子要分清两件事：

\begin{itemize}
\item \textbf{一般规律}：群变大只会把轨道并粗，绝不会把轨道分细；所以轨道个数单调不增。
\item \textbf{本例的巧合}：$D_8$ 相比 $C_4$ 没有真的把任何两条轨道并起来，所以个数没变。
\end{itemize}


想看“群变大真的把轨道并起来”，只要把黑白两色换成三色。仍是正方形四个顶点，$C_4$ 作用在 $3^4=81$ 种染色上得到 $24$ 条轨道，$D_8$ 得到 $21$ 条：反射确实把 $C_4$ 的若干轨道两两合并了。

\textbf{所以 Burnside 引理算出“个数相同”并不能说明“结构相同”；要判断作用是否一样，还得看轨道本身。}

\subsection{$S_n$ 中共轭的循环结构}


令 $S_n$ 作用在自身上，作用仍是共轭。对任意 $g\in S_n$，有

\[
 g(a_1\ a_2\ \cdots\ a_k)g^{-1}
=(g(a_1)\ g(a_2)\ \cdots\ g(a_k)).
\]

证明只需追踪一个点：右边把 $g(a_i)$ 送到 $g(a_{i+1})$，并把 $g(a_k)$ 送回 $g(a_1)$。因此共轭只会把循环中的标签整体重命名，不会改变循环长度；不交循环分解的长度多重集（包括固定点对应的 $1$-循环）称为\textbf{循环类型}。

反过来，若两个置换具有相同循环类型，就可以在每个对应循环之间定义一个双射 $g$，把第一个置换的每个循环逐项送到第二个置换的对应循环。于是上面的公式给出共轭。因此

\[
\sigma,\tau\in S_n\text{ 共轭}
\quad\Longleftrightarrow\quad
\sigma,\tau\text{ 具有相同的循环类型}.
\]

这也解释了为什么 $(123)$ 不可能与 $(12)(34)$ 共轭：前者的循环长度为 $3,1$，后者为 $2,2$。

\subsection{例：$S_4$ 的五个共轭类}


$S_4$ 的元素按循环类型分为五类。循环类型写成“只记录每个长度的循环有几个”的简写：$2\,1^2$ 读作“一个 2-循环、两个 1-循环”，$2^2$ 读作“两个 2-循环”，$1^4$ 读作“四个 1-循环”（即恒等元）。类大小可由“选择支撑集，再排列成循环”的方法计算；总和应回到 $|S_4|=24$。下面用 $\binom nk=\frac{n!}{k!(n-k)!}$ 表示从 $n$ 个元素中选 $k$ 个的组合数。

\begin{itemize}
\item $1^4$：恒等元 $e$，大小 $1$。
\item $2\,1^2$：一个换位，如 $(12)$，大小 $\binom42=6$。
\item $3\,1$：一个三循环，如 $(123)$，大小 $\binom43(3-1)!=8$。
\item $4$：一个四循环，如 $(1234)$，大小 $(4-1)!=6$。
\item $2^2$：两个不相交的换位，如 $(12)(34)$，大小 $\frac{4!}{2^2\,2!}=3$。
\end{itemize}


于是 $S_4$ 的类方程（按共轭类选代表元）是

\[
24=1+6+8+6+3.
\]

只有恒等元的共轭类大小为 $1$，所以 $S_4$ 的\textbf{中心}（与所有元素交换的元素全体，记作 $Z(G)$，第 3 节会正式定义）只有恒等元：

\[
Z(S_4)=\{e\}.
\]

更一般地，若 $\sigma\in S_n$ 的循环类型中长度 $i$ 的循环有 $m_i$ 个，则其共轭类大小为

\[
\frac{n!}{\prod_{i\ge1} i^{m_i}m_i!}.
\]

分母中的 $i^{m_i}$ 来自每个 $i$-循环的循环表示有 $i$ 种起点，$m_i!$ 来自同长度循环之间可以互换顺序。这个公式在手算类大小时很有用，但最重要的判别仍是“循环类型相同当且仅当共轭”。

\subsection{记忆点}


\begin{itemize}
\item Burnside：先算每个 $g$ 的固定点数，再平均；不要把固定点集和稳定子混为一谈。
\item $S_n$ 的共轭只改标签，不改不交循环的长度。
\item $S_4$ 的五种循环类型给出 $1+6+8+6+3=24$；这也是检查类大小计算是否漏项的快速方法。
\end{itemize}


\subsection{参考对照}


\begin{itemize}
\item [DF] 关于群作用、轨道计数、Burnside 引理与对称群共轭类的章节。
\item [DS] 关于群作用、轨道计数和置换循环分解的章节。
\end{itemize}

\section{共轭作用与类方程}


类方程来自一个非常具体的群作用：让 $G$ 通过共轭作用在自己身上。它把“一个元素能被共轭搬到哪些元素”与中心、中心化子联系起来，是证明有限 $p$-群结论和 Sylow 定理的入口。

上一节已经把数轨道的一般工具（Burnside 引理）准备好；这一节换成共轭作用，把群本身的元素按“能否互相共轭”分块，得到的计数恒等式就是类方程。

\subsection{群作用回顾}


左群作用是一个映射

\[
G\times X\longrightarrow X,\qquad (g,x)\longmapsto g\cdot x,
\]

满足 $e\cdot x=x$ 以及 $(gh)\cdot x=g\cdot(h\cdot x)$。对固定 $x$，定义轨道和稳定子：

\[
G\cdot x=\{g\cdot x:g\in G\},
\qquad
G_x=\{g\in G:g\cdot x=x\}.
\]

轨道是等价类，因此把 $X$ 划分成互不相交的轨道。轨道--稳定子定理给出集合双射 $G/G_x\cong G\cdot x$；当 $G$ 有限时

\[
|G\cdot x|=[G:G_x]=\frac{|G|}{|G_x|}.
\]

\subsection{共轭作用、共轭类与中心化子}


令 $X=G$，定义

\[
 g\cdot x=gxg^{-1}.
\]

验证作用公理只需计算

\[
(gh)x(gh)^{-1}=g(hxh^{-1})g^{-1}.
\]

元素 $x$ 的轨道

\[
\operatorname{Cl}_G(x)=\{gxg^{-1}:g\in G\}
\]

叫做 $x$ 的共轭类。它的稳定子为

\[
C_G(x)=\{g\in G:gxg^{-1}=x\}
=\{g\in G:gx=xg\},
\]

称为 $x$ 的中心化子（centralizer）。因此

\[
|\operatorname{Cl}_G(x)|=[G:C_G(x)].
\]

中心是所有共轭类只有一个元素的点：

\[
Z(G)=\{x\in G:gx=xg\ \text{对所有 }g\in G\}.
\]

所以共轭类大小为 $1$ 的元素恰好是 $Z(G)$ 的元素。

群也作用在子群集合（子群格）上：

\[
 g\cdot H=gHg^{-1}.
\]

这时 $H$ 的稳定子不是中心化子，而是正规化子

\[
N_G(H)=\{g\in G:gHg^{-1}=H\}.
\]

因此要区分：$C_G(x)$ 固定一个元素，$N_G(H)$ 固定一个子群作为集合。

\subsection{类方程的推导}


设 $x_1,\ldots,x_r$ 是所有非中心共轭类的代表元。共轭类把 $G$ 分割成中心中的单元素类和这些非中心类，于是

\[
|G|=|Z(G)|+\sum_{i=1}^{r}|\operatorname{Cl}_G(x_i)|
=|Z(G)|+\sum_{i=1}^{r}[G:C_G(x_i)].
\]

这就是 \textbf{类方程（class equation）}。它不是凭空写出的公式：第一项来自轨道大小为 $1$ 的元素，第二项来自其余轨道，再应用轨道--稳定子定理。

\subsection{有限 $p$-群的中心非平凡}


设 $|G|=p^k$，其中 $p$ 是素数、$k\ge 1$。每个中心化子 $C_G(x_i)$ 都是 $G$ 的子群，所以 $[G:C_G(x_i)]$ 是 $p$ 的正幂；因为 $x_i$ 非中心，该指数大于 $1$，因而被 $p$ 整除。类方程模 $p$ 化为

\[
|G|\equiv |Z(G)|\pmod p.
\]

左边为 $0$，故 $p\mid |Z(G)|$。尤其 $|Z(G)|\ge p$，所以中心不可能只有单位元。

\subsection{例：阶为 $p^2$ 的群必为阿贝尔群}


设 $|G|=p^2$。由上面的结论，$|Z(G)|$ 为 $p$ 或 $p^2$。

若 $|Z(G)|=p^2$，显然 $Z(G)=G$，所以 $G$ 阿贝尔。假设 $|Z(G)|=p$，取 $x\notin Z(G)$。中心化子 $C_G(x)$ 包含 $Z(G)$ 和 $x$，因此它严格大于 $Z(G)$。由拉格朗日定理，$|C_G(x)|$ 必须是 $p^2$，从而 $C_G(x)=G$。这意味着 $x$ 与每个群元素交换，即 $x\in Z(G)$，矛盾。

所以 $|Z(G)|\ne p$，只能是 $p^2$，进而 $G$ 阿贝尔。结合有限生成阿贝尔群结构定理，阶为 $p^2$ 的群只有

\[
C_{p^2}\quad\text{或}\quad C_p\times C_p
\]

两种同构类型。

\subsection{记忆点}


\begin{itemize}
\item 共轭作用的轨道是共轭类，稳定子是中心化子：$|\operatorname{Cl}_G(x)|=[G:C_G(x)]$。
\item 中心是共轭类大小为 $1$ 的元素集合。
\item 子群的共轭作用稳定子是正规化子 $N_G(H)$，不要与 $C_G(x)$ 混用。
\item 类方程只对每个非中心共轭类选一个代表元；不能把同一共轭类的所有元素重复相加。
\end{itemize}


\subsection{参考对照}


\begin{itemize}
\item [DF] 关于群作用、共轭类、中心化子与类方程的章节。
\item [DS] 关于群作用与类方程的章节；阶为 $p^2$ 的结论是类方程的直接应用。
\end{itemize}

\section{Sylow 子群的存在性}


\subsection{先固定记号}


设 $p$ 是素数，有限群 $G$ 的阶写成

\[
|G|=p^a m,\qquad p\nmid m.
\]

若 $a\ge 1$，阶为 $p^a$ 的子群称为 \textbf{Sylow $p$-子群}；记所有 Sylow $p$-子群组成的集合为
$\operatorname{Syl}_p(G)$，其个数记作 $n_p$。这里的 $p^a$ 是 $|G|$ 中能出现的最大 $p$-幂，因此 Sylow 子群是“尽可能大的 $p$-子群”，而不是任意一个阶为 $p$ 的子群。

当 $a=0$ 时，$p\nmid |G|$，只有平凡子群 $\{e\}$，把它视为唯一的 Sylow $p$-子群即可。下面只讨论 $a\ge1$。

\subsection{定理：Sylow 子群存在}


\textbf{Sylow 第一定理（存在性部分）}：若 $|G|=p^a m$ 且 $p\nmid m$，则 $G$ 至少有一个阶为 $p^a$ 的子群。

证明采用 $|G|$ 的归纳。需要用到第 3 节的类方程

\[
|G|=|Z(G)|+\sum_i [G:C_G(x_i)],
\]

其中 $x_i$ 取所有非中心共轭类的代表，$C_G(x)=\{g\in G:gx=xg\}$ 是元素 $x$ 的中心化子。证明的关键是：要么中心中能取出一个阶为 $p$ 的正规子群，从而降到商群；要么类方程迫使某个较小的中心化子保留完整的 $p^a$ 因子。

\subsection{情形一：$p\mid |Z(G)|$}


中心 $Z(G)=\{x\in G:gx=xg\text{ 对所有 }g\in G\}$ 是有限阿贝尔群。由\textbf{阿贝尔群版本的 Cauchy 定理}（若素数 $p$ 整除有限阿贝尔群 $A$ 的阶，则 $A$ 中存在阶为 $p$ 的元素），$Z(G)$ 中存在阶为 $p$ 的元素 $z$。令

\[
N=\langle z\rangle=\{e,z,z^2,\dots,z^{p-1}\},\qquad |N|=p,
\]

其中 $\langle z\rangle$ 表示由 $z$ 生成的循环子群。

因为 $N\le Z(G)$，所以 $N\trianglelefteq G$，且

\[
|G/N|=\frac{|G|}{|N|}=p^{a-1}m.
\]

商群比 $G$ 小，归纳假设给出 $G/N$ 的一个 Sylow $p$-子群 $\overline P$，满足 $|\overline P|=p^{a-1}$。令自然映射 $\pi:G\to G/N$，取其原像

\[
P=\pi^{-1}(\overline P).
\]

由同态基本计数公式（或原像的阶公式）

\[
|P|=|N|\,|\overline P|=p\cdot p^{a-1}=p^a.
\]

于是 $P$ 就是 $G$ 的 Sylow $p$-子群。

\textbf{为什么这里要取原像？} 归纳只在商群 $G/N$ 中给出了一个子群；原像把它“抬回”到 $G$，并自动包含核 $N$，所以阶数多出恰好一个因子 $p$。

\subsection{情形二：$p\nmid |Z(G)|$}


（本节反复出现 $[G:H]=|G|/|H|$，即子群 $H\le G$ 的\textbf{指数}；$p\nmid n$ 读作“$p$ 不整除 $n$”。）

此时使用类方程。若对每个非中心代表 $x_i$ 都有

\[
p\mid [G:C_G(x_i)],
\]

那么右端的每一项都被 $p$ 整除，因而 $p\mid |Z(G)|$；这与当前情形矛盾。因此存在某个非中心元素 $x$，使得

\[
p\nmid [G:C_G(x)].
\]

记 $H=C_G(x)$。由于 $x\notin Z(G)$，$H$ 是 $G$ 的真子群，所以可以对 $H$ 使用归纳假设。又因为

\[
|G|=|H|[G:H],\qquad p\nmid [G:H],
\]

而 $|G|$ 的 $p$-部分是 $p^a$，可知 $|H|$ 仍含有完整的 $p^a$ 因子：存在 $k$ 使

\[
|H|=p^a k,\qquad p\nmid k.
\]

归纳假设给出 $H$ 中一个阶为 $p^a$ 的子群 $P$。把 $P$ 看成 $G$ 的子群，它的阶已经等于 $|G|$ 的最大 $p$-幂，因此 $P$ 是 $G$ 的 Sylow $p$-子群。

两种情形穷尽了所有可能性，存在性定理得证。

\subsection{课堂证明中容易错的地方}


\begin{itemize}
\item 归纳的第二种情形中，$H=C_G(x)$ 的阶是 $p^a k$，不是只剩下 $p^{a-1}$；因为选择的是一个不被 $p$ 整除的共轭类大小。
\item $C_G(x)$ 是\textbf{中心化子}，由“与元素 $x$ 交换”定义；后面对 Sylow 子群集合做共轭作用时出现的是\textbf{正规化子} $N_G(P)$，两者不能混用。
\item $G/N$ 的阶是 $p^{a-1}m$，不要把其中的 $m$ 漏掉。归纳得到的是商群中的阶 $p^{a-1}$ 子群，再取原像才得到 $G$ 中的阶 $p^a$ 子群。
\end{itemize}


\subsection{例：$S_4$ 中的 Sylow 子群阶数}


\[
|S_4|=24=2^3\cdot3.
\]

Sylow $2$-子群的阶必须是 $8$，Sylow $3$-子群的阶必须是 $3$。这里 $S_4$ 是 $\{1,2,3,4\}$ 上的对称群，轮换记号 $(1\,2\,3)$ 表示 $1\mapsto2\mapsto3\mapsto1$。例如

\[
\langle(1\,2\,3)\rangle=\{e,(1\,2\,3),(1\,3\,2)\}
\]

是一个 Sylow $3$-子群。存在性定理保证阶为 $8$ 的 Sylow $2$-子群也一定存在；至于这样的子群有多少个、彼此是否共轭，要等第二、第三定理处理。

\subsection{课堂补充：$S_n$ 的 Sylow $p$-子群怎么构造}


存在性定理说 Sylow 子群一定存在，但没说怎么找。在对称群里有一个可以手写的答案：\textbf{让置换只在自己的 $p$ 元块内部活动，再让它在块之间整体轮换}。课堂讲的就是这个构造，下面把它写成可复用的形式。

先做最简单的情形 $n=p^k$（$S_{p^k}$ 是 $p^k$ 个字母上的对称群），此时 $|S_{p^k}|$ 中 $p$ 的幂次是

\[
p^{1+p+p^2+\cdots+p^{k-1}}.
\]

把 $p^k$ 个点分成 $p$ 个块 $B_1,\ldots,B_p$，每块 $p^{k-1}$ 个点。递归地定义 $W_k\le S_{p^k}$：

\begin{itemize}
\item $k=1$：$W_1=\langle(1\,2\,\cdots\,p)\rangle$ 是 $p$ 阶循环群；
\item 从 $W_{k-1}$ 到 $W_k$：把 $W_{k-1}$ 的\textbf{每个生成元在每块内部同时作用}，再加上一个\textbf{把 $p$ 个块整体轮换}的置换。
\end{itemize}


这样得到的 $W_k$ 的阶是 $p^{k-1}\cdot|W_{k-1}|$，逐层展开正好是 $p^{1+p+\cdots+p^{k-1}}$，也就是 $S_{p^k}$ 的 Sylow $p$-子群。

\textbf{$S_4$，$p=2$。} $k=2$，$|S_4|$ 的 $2$-部分是 $2^3=8$。取两块 $\{1,2\}$、$\{3,4\}$：


\begin{center}
\begin{tabular}{ll}
生成元 & 含义\\ \hline
$(1\,2)$ & 第 1 块内部交换\\
$(3\,4)$ & 第 2 块内部交换\\
$(1\,3)(2\,4)$ & 两块整体轮换\\
$(1\,4)(2\,3)$ & 两块整体反向轮换
\end{tabular}
\end{center}


它们生成的子群阶为 $8$（直接乘出 $8$ 个元素即可验证），正是 $S_4$ 的 Sylow $2$-子群。注意这里 $4$ 个生成元比需要的多：$(1\,4)(2\,3)$ 可以由前三个乘出来，列出来只是为了看清“块内 $\times$ 块间”的结构。

\textbf{$S_6$，$p=3$（一个容易踩的坑）。} $6=3\cdot2$，$|S_6|$ 的 $3$-部分是 $3^2=9$，不是 $3\cdot3\cdot2=18$。如果天真地把“每块内部的 $3$-轮换”和“交换两块”都当成独立生成元，会得到一个 $18$ 阶子群，\textbf{它不是 $S_6$ 的 Sylow $3$-子群}（阶太大了）。正确的做法是让两块内部的 $3$-轮换相互关联，使整个子群只有 $9$ 个元素。这说明构造 Sylow 子群时，“阶必须整除 $|G|$ 的最大 $p$-幂”是必须回头检查的约束。

\subsection{参考对照}


\begin{itemize}
\item [DF] §4.5，Theorem 18（Sylow 定理及其归纳证明）；$S_n$ 的 Sylow 子群可用 $p$ 元块上的迭代\textbf{圈积}（wreath product，记作 $\wr$）描述——本节写出的递归构造就是它的具体展开，不需要预先知道圈积的定义。
\item [DS] 群作用与轨道计数的相关小节。该书正文没有 Sylow 定理，本节沿用课堂的类方程路线，并补全商群原像的阶数计算；$S_n$ 的构造按课堂补充整理。
\end{itemize}

\section{Sylow 子群的共轭性}


\textbf{本节用到的记号，先一次交代清楚}：

\begin{itemize}
\item $G$：有限群；$P\in\operatorname{Syl}_p(G)$：一个 Sylow $p$-子群；$Q\le G$：任意 $p$-子群。
\item $\operatorname{Syl}_p(G)$：全体 Sylow $p$-子群；$G$ 通过共轭作用在它上面，$g\cdot P=gPg^{-1}$。
\item $N_G(P)=\{g\in G:gPg^{-1}=P\}$：$P$ 的正规化子，它正是上述作用中 $P$ 的稳定子。
\item $C_G(P)$：$P$ 的中心化子。注意：共轭作用的稳定子是 $N_G(P)$ 而\textbf{不是} $C_G(P)$；只有 $P$ 是 Abel 群时二者才相同。
\item $|G|=p^am,\ p\nmid m$：阶的分解；$[G:N_G(P)]$：指数。
\item $n_p=|\operatorname{Syl}_p(G)|$：Sylow $p$-子群的个数。
\item $(c_1\ c_2\ \cdots\ c_k)$：轮换记号；$S_n$：对称群；$\trianglelefteq$：正规子群；$\le$：子群。
\end{itemize}


\subsection{第二定理的内容}


设 $|G|=p^a m$，其中 $p\nmid m$，并令 $P\in\operatorname{Syl}_p(G)$。Sylow 第二定理包含两句话：

1. 任意 $p$-子群 $Q\le G$ 都包含在某个 $P$ 的共轭子群中，即存在 $g\in G$ 使

\[
   Q\le gPg^{-1}.
\]

2. 任意两个 Sylow $p$-子群彼此共轭。特别地，它们同构，但通常不是同一个子群。

第一句话更强：它告诉我们任意 $p$-子群都可以“放大”到某个 Sylow 子群里；第二句话是取 $|Q|=p^a$ 得到的。

\subsection{用左陪集证明包含结论}


令 $Q\le G$ 是任意 $p$-子群。考虑左陪集集合

\[
\Omega=G/P=\{gP:g\in G\}.
\]

这里 $G/P$ 只是左陪集的集合，\textbf{暂时不能当作商群}，因为 $P$ 未知是否正规。陪集个数为

\[
|\Omega|=[G:P]=m,
\]

不被 $p$ 整除。$Q$ 通过左乘作用在 $\Omega$ 上：

\[
q\cdot(gP)=(qg)P.
\]

这个确实是群作用，因为左乘保持陪集的形式。由轨道--稳定子定理，每个 $Q$-轨道的大小都是 $|Q|$ 的因子，因而是 $p$ 的幂。把 $\Omega$ 分解为不相交轨道后，若所有轨道大小都被 $p$ 整除，则 $|\Omega|$ 也被 $p$ 整除，矛盾。因此存在一个大小不被 $p$ 整除的轨道；由于它的大小是 $p$ 的幂，只能是 $1$。

所以有某个 $gP\in\Omega$ 被 $Q$ 的每个元素固定：

\[
qgP=gP\qquad(\forall q\in Q).
\]

左乘 $g^{-1}$，得到

\[
g^{-1}qgP=P.
\]

一个左乘 $hP$ 等于 $P$ 当且仅当 $h\in P$，故 $g^{-1}qg\in P$。对所有 $q\in Q$ 成立，于是

\[
g^{-1}Qg\le P,
\qquad\text{即}\qquad
Q\le gPg^{-1}.
\]

这就证明了：每个 $p$-子群都包含于某个 Sylow $p$-子群的共轭中。

\subsubsection{等价的“交”形式（[DF] 的 Lemma 19）}


上面的结论有一个看起来不同、其实等价的写法，它在 [DF] 里是证明 Sylow 定理的起点：

\textbf{引理。} 设 $P\in\operatorname{Syl}_p(G)$，$Q\le G$ 是任意 $p$-子群，则

\[
Q\cap N_G(P)=Q\cap P.
\]

\textbf{证明。} 因为 $P\le N_G(P)$，所以 $P\cap Q\le N_G(P)\cap Q$，只需证反向包含。取 $x\in Q\cap N_G(P)$。由 $x\in N_G(P)$ 知 $xPx^{-1}=P$，于是 $xPx^{-1}\le P$；而 $x\in Q$，所以 $xPx^{-1}\le Q$。两边合起来

\[
xPx^{-1}\le P\cap Q.
\]

取阶得 $|P|\le|P\cap Q|$，结合 $P\cap Q\le P$ 得到 $P\cap Q=P$，即 $P\le Q$。又 $x\in Q$ 且 $P\le Q$，所以 $x\in P$。这就证明了 $Q\cap N_G(P)\le Q\cap P$。$\square$

两种写法的关系是直接的：$Q\le gPg^{-1}$ 等价于 $g^{-1}Qg\le P$，也就是“某个共轭的 $Q$ 被 $P$ 吸收”；而引理说的是“$Q$ 中那些把 $P$ 保持为集合的元素，本来就已经在 $P$ 里”。读 [DF] 时如果对 Lemma 19 的作用感到突然，可以回到上面陪集作用的版本：两者证的是同一件事。

\subsection{两个 Sylow 子群必共轭}


若 $Q$ 本身也是 Sylow $p$-子群，则 $|Q|=|P|=p^a$。上面的包含关系给出

\[
Q\le gPg^{-1},\qquad |Q|=|gPg^{-1}|.
\]

有限群中阶相同的包含子群只能相等，所以 $Q=gPg^{-1}$。因此

\[
\operatorname{Syl}_p(G)=\{gPg^{-1}:g\in G\},
\]

是 $G$ 在共轭作用下的一个轨道。

\subsection{正规化子：不要和中心化子混淆}


对子群 $P\le G$，定义正规化子

\[
N_G(P)=\{g\in G:gPg^{-1}=P\}.
\]

它是 $G$ 的子群，并且 $P\le N_G(P)$。由定义，$P$ 在 $N_G(P)$ 中正规：

\[
P\trianglelefteq N_G(P).
\]

这里有两个相似但不同的对象：

\begin{itemize}
\item 元素 $x$ 的中心化子 $C_G(x)=\{g:gx=xg\}$，要求每个元素逐点固定 $x$；
\item 子群 $P$ 的正规化子 $N_G(P)=\{g:gPg^{-1}=P\}$，只要求整体保持子群 $P$ 不变。
\end{itemize}


在 $G$ 共轭作用于 $\operatorname{Syl}_p(G)$ 时，$P$ 的稳定子正是 $N_G(P)$，不是 $C_G(P)$。课堂转写中“除以中心化子”的地方应改为“除以正规化子”。

\subsection{唯一性与正规性}


以下三个条件等价：

\[
n_p=1
\quad\Longleftrightarrow\quad
P\trianglelefteq G
\quad\Longleftrightarrow\quad
gPg^{-1}=P\ (\forall g\in G).
\]

证明很短。若 $n_p=1$，那么 $gPg^{-1}$ 仍是 Sylow $p$-子群，只能等于唯一的 $P$，所以 $P$ 正规。反过来若 $P$ 正规，则所有共轭 $gPg^{-1}$ 都等于 $P$；而第二定理说所有 Sylow 子群都来自这些共轭，因此只有一个。

\textbf{“唯一”指的是什么？} 这里的唯一是集合意义上的“只有一个子群”。一般情况下不同 Sylow 子群虽不相等，却都共轭；说“唯一到共轭”则是说同一个素数对应的 Sylow 子群只有一个共轭类，这个结论在第二定理中总是成立。

\subsection{例：$S_4$ 的 Sylow $2$-子群}


\[
|S_4|=24=2^3\cdot3,
\]

所以 Sylow $2$-子群阶为 $8$。一个具体构造是把两个二元块 $\{1,2\}$ 与 $\{3,4\}$ 互换：块内的交换 $(1\,2)$、$(3\,4)$，以及交换两块的置换 $(1\,3)(2\,4)$ 共同生成一个阶为 $8$ 的子群。任意另一个阶为 $8$ 的子群都与它共轭，但未必等于这个具体子群。

\subsection{参考对照}


\begin{itemize}
\item [DF] §4.5，Theorem 18(2) 及其陪集/轨道证明。
\item [DS] 该书没有正规化子的讨论；正规化子的定义与 "$P\trianglelefteq N_G(P)$" 来自 [DF] §4.4 及本节给出的定义。
\end{itemize}

\section{Sylow 子群的个数}


\subsection{第三定理}


继续设

\[
|G|=p^a m,\qquad p\nmid m,
\]

并记 $n_p=|\operatorname{Syl}_p(G)|$。Sylow 第三定理给出

\[
n_p\mid m,
\qquad
n_p\equiv1\pmod p.
\]

这两个条件经常一起使用：先把 $m$ 的因数列出来，再筛掉不满足模 $p$ 条件的候选值。

\subsection{先证明 $n_p\mid m$}


固定一个 Sylow 子群 $P$。由第二定理，$G$ 通过共轭作用在 $\operatorname{Syl}_p(G)$ 上只有一条轨道，这时称作用是\textbf{传递的}（transitive）；$P$ 的稳定子是

\[
\operatorname{Stab}_G(P)=\{g:gPg^{-1}=P\}=N_G(P).
\]

轨道--稳定子定理于是给出

\[
n_p=[G:N_G(P)].
\]

因为 $P\le N_G(P)\le G$，指数可以分解为

\[
[G:P]=[G:N_G(P)]\,[N_G(P):P].
\]

左边 $[G:P]=m$，所以

\[
m=n_p\,[N_G(P):P],
\]

从而 $n_p\mid m$。

这里再次出现的是正规化子 $N_G(P)$。不能把稳定子写成中心化子 $C_G(P)$：中心化子要求逐个元素与 $P$ 中所有元素交换，而共轭作用只要求把整个子群 $P$ 保持为集合。

\subsection{再证明 $n_p\equiv1\pmod p$}


让 $P$ 自己通过共轭作用在集合

\[
\mathcal S=\operatorname{Syl}_p(G)
\]

上。把 $\mathcal S$ 分解成 $P$-轨道。轨道 $\mathcal O(Q)$ 的大小为

\[
|\mathcal O(Q)|=[P:\operatorname{Stab}_P(Q)]
=[P:P\cap N_G(Q)].
\]

由于 $P$ 是 $p$-群，每个轨道大小都是 $p$ 的幂。我们要证明：除了 $Q=P$ 这一点轨道外，其余轨道大小都被 $p$ 整除。

\subsubsection{哪些轨道可能是单点？}


显然 $P$ 的共轭轨道是单点，因为对任意 $x\in P$，$xPx^{-1}=P$；也就是说 $P\le N_G(P)$。

反过来，假设某个 $Q\in\mathcal S$ 的轨道也是单点。那么

\[
P\le N_G(Q).
\]

子群 $Q$ 按定义正规于 $N_G(Q)$。因此 $P,Q$ 都是 $N_G(Q)$ 的 $p$-子群，并且 $P$ 已经有阶 $p^a$。由于 $N_G(Q)\le G$，它的阶不可能含有比 $p^a$ 更高的 $p$-幂，所以 $P$ 和 $Q$ 都是 $N_G(Q)$ 的 Sylow $p$-子群。由第二定理，它们在 $N_G(Q)$ 中共轭；但 $Q\trianglelefteq N_G(Q)$，其共轭只能还是 $Q$，故 $P=Q$。

因此

\[
\text{单点轨道只有 }\{P\}.
\]

对任意 $Q\ne P$，其轨道不是单点，所以轨道大小这个 $p$ 的幂至少为 $p$，即被 $p$ 整除。把所有轨道大小相加：

\[
n_p=|\mathcal S|
=1+\sum_{\mathcal O\ne\{P\}}|\mathcal O|
\equiv1\pmod p.
\]

这就得到第二个结论。

\subsection{数数时的标准流程}


面对“求 $n_p$ 的可能值”这类题，可以固定以下顺序：

1. 分解 $|G|=p^a m$，先确认 Sylow $p$-子群的阶是 $p^a$。
2. 写出 $n_p\mid m$ 的所有正因数候选。
3. 筛选满足 $n_p\equiv1\pmod p$ 的候选。
4. 若只剩 $n_p=1$，立刻得到 Sylow $p$-子群正规；若仍有多个候选，再结合群作用、共轭类或元素计数。

\subsection{例：$S_3$}


\[
|S_3|=6=2\cdot3.
\]

这里 $S_3$ 是 $\{1,2,3\}$ 上的对称群，$A_3$ 是其中的偶置换子群。对 $p=3$，有 $n_3\mid2$ 且 $n_3\equiv1\pmod3$，只能是 $n_3=1$。因此唯一的 Sylow $3$-子群

\[
\langle(1\,2\,3)\rangle=A_3
\]

正规于 $S_3$。

对 $p=2$，有 $n_2\mid3$ 且 $n_2\equiv1\pmod2$，候选为 $1$ 或 $3$。实际上三个换位分别生成三个不同的阶 $2$ 子群，所以 $n_2=3$；它们彼此共轭而没有一个正规。

\subsection{例：阶为 $35$ 的群}


设 $|G|=35=5\cdot7$。对 $p=7$，$n_7\mid5$ 且 $n_7\equiv1\pmod7$，故 $n_7=1$。对 $p=5$，$n_5\mid7$ 且 $n_5\equiv1\pmod5$，也只能有 $n_5=1$。

记 $P_5$ 为那个唯一的 Sylow $5$-子群、$P_7$ 为唯一的 Sylow $7$-子群（下标就是对应的素数 $p$，与第 7 节 $P_{11},P_{13}$ 的写法一致）。两者都正规；它们阶\textbf{互素}（即 $\gcd(5,7)=1$），所以交集只能是 $\{e\}$；又因为两个正规子群的元素彼此交换，乘积是直积，于是

\[
G\cong P_5\times P_7\cong C_5\times C_7\cong C_{35}.
\]

最后一步用了 $C_m\times C_n\cong C_{mn}$（当 $\gcd(m,n)=1$）。

这里“先由数目得到正规，再用正规子群的乘积”是 Sylow 定理最典型的应用。更大的例子（例如 $11^2\cdot13^2$）还要结合“阶为 $p^2$ 的群必阿贝尔”等结果，不能只凭 $n_p$ 的两个条件直接断言循环。

\subsection{参考对照}


\begin{itemize}
\item [DF] §4.5，Theorem 18(3) 及其证明。
\item [DS] Sylow 定理与小阶群应用相关小节。
\end{itemize}

\section{Sylow 定理的应用：小阶群与 $pq$ 群}


\textbf{本节用到的记号，先一次交代清楚}：

\begin{itemize}
\item $G$：有限群；$|G|$：阶；$H,K\le G$：子群；$P,Q$：Sylow 子群（分别对应两个不同的素数）。
\item $n_p$：Sylow $p$-子群的个数；判据 $n_p\mid m$ 且 $n_p\equiv1\pmod p$；$n_p=1\iff$ 该 Sylow 子群正规。
\item $P_p$：当 $n_p=1$ 时记那个唯一的 Sylow $p$-子群（如 $P_5,P_7,P_{11},P_{13}$）。
\item $HK=\{hk\}$：子群乘积；$H\cap K$：交；$[x,y]=xyx^{-1}y^{-1}$：交换子，$[x,y]=e\iff xy=yx$。
\item 内部直积：$HK=G$、$H\cap K=\{e\}$（且两者都正规）时 $G\cong H\times K$。
\item $C_n$：$n$ 阶循环群；$C_q\rtimes_\varphi C_p$：半直积（$\varphi$ 平凡时退化为直积 $\times$）。
\item $\varphi:C_p\to\operatorname{Aut}(C_q)$：把共轭作用写成自同构；$s\in(\mathbb Z/q\mathbb Z)^\times$ 满足 $s^p\equiv1\pmod q$。
\item $\operatorname{Aut}(C_q)\cong(\mathbb Z/q\mathbb Z)^\times$：阶为 $q-1$ 的群；$\cong$：同构；$\trianglelefteq$：正规子群。
\end{itemize}


这一节把 Sylow III 的两个条件
\[
n_p\mid m,\qquad n_p\equiv 1\pmod p
\]
当作“筛选器”使用。先列出可能的 $n_p$，再利用正规性和共轭作用完成结构判断。做题时不要一上来猜群的乘法表。

\subsection{阶为 $35$ 的群}


设 $|G|=35=5\cdot 7$。对 Sylow 子群分别计数：

\begin{itemize}
\item $n_7\mid 5$ 且 $n_7\equiv1\pmod7$，所以 $n_7=1$；
\item $n_5\mid 7$ 且 $n_5\equiv1\pmod5$，所以 $n_5=1$。
\end{itemize}


因此 Sylow $7$-子群 $Q$ 与 Sylow $5$-子群 $P$ 都正规。它们的交只有单位元，因为 $|P\cap Q|$ 同时整除 $5$ 和 $7$。对任意 $x\in P,y\in Q$，交换子
\[
[x,y]=xyx^{-1}y^{-1}
\]
属于 $P$（因为 $P\trianglelefteq G$），也属于 $Q$（因为 $Q\trianglelefteq G$），所以 $[x,y]=e$。两子群逐点交换，于是 $PQ$ 是子群，而且
\[
|PQ|=\frac{|P||Q|}{|P\cap Q|}=35=|G|,
\]
从而 $G=PQ\cong P\times Q\cong C_5\times C_7\cong C_{35}$。这说明阶 $35$ 的群必为循环群。

\subsection{阶为 $11^2\cdot13^2$ 的群：先判定阿贝尔，不急着判定循环}


设 $|G|=11^2\cdot13^2=20449$。Sylow 条件给出
\[
 n_{11}\mid13^2,\quad n_{11}\equiv1\pmod{11},
\]
所以 $n_{11}\in\{1,13,169\}$ 中只有 $1$ 满足同余条件。类似地，
\[
 n_{13}\mid11^2,\quad n_{13}\equiv1\pmod{13},
\]
而 $n_{13}\in\{1,11,121\}$ 中只有 $1$ 满足条件。因此两个 Sylow 子群都正规；下面记 $P_p$ 为那个唯一的 Sylow $p$-子群（即 $P_{11}$ 是 Sylow $11$-子群、$P_{13}$ 是 Sylow $13$-子群，下标就是素数 $p$）。

阶为 $p^2$ 的群是阿贝尔群（见本讲类方程部分），所以 $P_{11}$ 与 $P_{13}$ 都阿贝尔。与阶 $35$ 的论证相同，两个正规、互素阶子群逐点交换，并且乘积等于 $G$。故 $G$ 阿贝尔。

这里要停在“阿贝尔”：$P_{11}$ 可以是 $C_{121}$ 或 $C_{11}\times C_{11}$，$P_{13}$ 也有两种选择，所以不能仅凭阶数断言 $G\cong C_{20449}$。这正是“先判定整体性质，再区分阿贝尔分类”的例子。

\subsection{阶为 $pq$ 的群（$p<q$ 为素数）}


设 $|G|=pq$。Sylow 条件先给出
\[
 n_q\mid p,\qquad n_q\equiv1\pmod q.
\]
因为 $p<q$，只有 $n_q=1$，所以 Sylow $q$-子群 $Q=\langle a\rangle$ 正规，且 $|a|=q$（$\langle a\rangle$ 表示由 $a$ 生成的循环子群）。令 $P=\langle b\rangle$ 为 Sylow $p$-子群，$|b|=p$。由于 $|P||Q|=|G|$ 且交只有单位元，$G=QP$，每个元素都可写成 $a^i b^j$（$i=0,\dots,q-1$，$j=0,\dots,p-1$）。

\subsubsection{共轭作用给出分类分岔}


正规性保证 $bab^{-1}\in Q$。$Q$ 是阶 $q$ 的循环群，因此存在整数 $s$ 使 $bab^{-1}=a^s$。这里 $s$ 只能取自模 $q$ 的\textbf{单位群}

\[
(\mathbb Z/q\mathbb Z)^\times=\{s\bmod q:\gcd(s,q)=1\},
\]

它由所有与 $q$ 互素的剩余类组成，在乘法下构成群，阶为 $q-1$（$q$ 为素数时每个非零剩余类都可逆）。取这样的 $s$（而不是任意的 $s$）是因为 $a\mapsto a^s$ 必须是 $Q$ 的自同构。于是
\[
bab^{-1}=a^s.
\]
因为 $b^p=e$，反复共轭得到
\[
a=b^pab^{-p}=a^{s^p},
\]
即 $s^p\equiv1\pmod q$。于是只有两种情形：

\begin{itemize}
\item $s=1$。此时 $a$ 与 $b$ 交换，$G\cong C_q\times C_p\cong C_{pq}$；
\item $s\ne1$。此时 $s$ 在单位群 $(\mathbb Z/q\mathbb Z)^\times$ 中的阶为 $p$（因为 $s^p\equiv1$ 而 $s\ne1$）。单位群的阶是 $q-1$，由拉格朗日定理 $p$ 整除 $q-1$，所以必须有 $p\mid(q-1)$。这给出非阿贝尔的\textbf{半直积}关系
\end{itemize}

\[
  \langle a,b\mid a^q=b^p=e,\;bab^{-1}=a^s\rangle.
\]

  （\textbf{半直积}的记号写在下面 [DF] 的写法里：$C_q\rtimes_\varphi C_p$。粗说就是“$C_q$ 正规、$C_p$ 通过自同构 $\varphi$ 作用于它”，当 $\varphi$ 不是恒等映射时结果不是直积。它的一般定义这里不展开，本讲只需要上面这一组生成元与关系。）

课堂上最后一个“唯一性”步骤只写到改变 $b$ 的生成元。这里用两种方式补齐，先给一个只用到多项式根的初等写法，再给教材的写法。

\textbf{初等写法。} 若 $s$ 与 $t$ 都是阶 $p$ 的非平凡元，则 $s$ 的 $p$ 个幂已经是方程 $X^p=1$ 在模 $q$ 下的全部根（一个次数为 $p$ 的多项式至多有 $p$ 个根）。所以 $t=s^r$，其中 $1\le r<p$。把生成元 $b$ 换成 $b^r$，共轭指数就变成 $s^r=t$。因此在固定的 $p,q$ 下，非阿贝尔情形至多有一个同构类型；它存在当且仅当 $p\mid(q-1)$。

\textbf{[DF] 的写法：换成自同构群。} 这里的 $\operatorname{Aut}(C_q)$ 表示 $C_q$ 的\textbf{自同构群}（$C_q$ 到自身的全体同构在复合下构成的群），同样有 $\operatorname{Aut}(C_q)\cong(\mathbb Z/q\mathbb Z)^\times$。正规性把共轭作用变成一个同态

\[
\varphi:C_p\longrightarrow\operatorname{Aut}(C_q),\qquad \varphi(b)=\bigl(a\mapsto bab^{-1}\bigr).
\]

因为 $C_q$ 是循环群，$\operatorname{Aut}(C_q)\cong(\mathbb Z/q\mathbb Z)^\times$ 是 $q-1$ 阶群。同态的像的阶同时整除 $p$ 与 $q-1$，所以只有两种可能：像平凡，或像的阶为 $p$。

\begin{itemize}
\item 像平凡：$b$ 与 $a$ 交换，群是直积 $C_p\times C_q\cong C_{pq}$；
\item 像的阶为 $p$：这要求 $p\mid(q-1)$，群是半直积 $C_q\rtimes_\varphi C_p$。
\end{itemize}


至于唯一性：$(\mathbb Z/q\mathbb Z)^\times$ 是循环群，所以其中阶为 $p$ 的子群至多一个，记作 $H$。任何像的阶为 $p$ 的同态，其像都等于这个唯一的 $H$；而 $H$ 循环，$\varphi(b)$ 必须是 $H$ 的生成元。于是两个这样的同态 $\varphi,\psi$ 满足 $\psi(b)=\varphi(b)^r$（某个 $1\le r<p$），即 $\psi=\varphi^r$。把 $b$ 换成 $b^r$ 就把 $\varphi$ 变成 $\psi$，两个半直积同构。\textbf{所以 $pq$ 阶群恰好有两种同构类型（循环群，以及 $p\mid(q-1)$ 时唯一的非阿贝尔群）}，这正是 [DF] §4.5 里“$p\nmid q-1$ 时必循环”的结论。

两种写法算的是同一件事：初等写法直接数“方程 $X^p=1$ 的根”，教材写法数“阶为 $p$ 的自同构”，而 $(\mathbb Z/q\mathbb Z)^\times$ 循环这件事实保证了这个“阶为 $p$ 的东西”只有一个。

为了看见这个群确实存在，可以在集合 $\mathbb Z/q\mathbb Z\times\mathbb Z/p\mathbb Z$ 上定义
\[
(i,j)(k,\ell)=(i+s^j k,\,j+\ell),
\]
其中第一坐标模 $q$、第二坐标模 $p$，并取 $s^p\equiv1\pmod q$。$\rtimes$ 读作“半直积”，$C_q\rtimes_\varphi C_p$ 中下标 $\varphi$ 指明 $C_p$ 通过哪个自同构作用在 $C_q$ 上。这正是 $C_q\rtimes C_p$ 的乘法；它有 $pq$ 个元素，且 $(1,0)$、$(0,1)$ 满足上面的关系。

\subsubsection{两个小例子}


\begin{itemize}
\item $|G|=6$：$p=2,q=3$，非平凡指数只能是 $s=2\equiv-1\pmod3$，得到
\end{itemize}

\[
  \langle a,b\mid a^3=b^2=e,\;bab^{-1}=a^{-1}\rangle\cong S_3.
\]
\begin{itemize}
\item $|G|=21$：$p=3,q=7$，因为 $3\mid6$，除了循环群外还可能有非阿贝尔群。取 $s=2$（$2^3\equiv1\pmod7$），关系为 $bab^{-1}=a^2$；取 $s=4=2^2$ 只是把 $b$ 换成 $b^2$ 后的同一同构类型。
\end{itemize}


\subsection{做题检查清单}


1. 写出 $|G|=p^a m$，先用 Sylow III 筛掉不可能的 $n_p$。
2. $n_p=1$ 才能直接写出 $P\trianglelefteq G$；“所有 Sylow 子群共轭”不等于“只有一个”。
3. 需要证明阿贝尔时，检查正规子群的交与交换子；需要分类时，再研究共轭作用 $P\to\operatorname{Aut}(Q)$。
4. 关系式中的指数在 $(\mathbb Z/q\mathbb Z)^\times$ 中计算，最后别忘了验证 $s^p\equiv1$。


\section*{参考文献}
\begin{itemize}
\item \textbf{[DF]} D. S. Dummit, R. M. Foote. \emph{Abstract Algebra}, 3rd ed. §§3.4, 4.1--4.5, 5.1, 5.5. John Wiley \& Sons, 2004.
\item \textbf{[DS]} 邓小铁、王畅。《离散数学与结构》。§§9.1--9.6, 14.1--14.2.
\end{itemize}
\end{document}
